%%%%%%%%%%%%%%%%%%%%%%%%%%%%%%%%%%%%%%%%%%%%%%%
% include additional packages you need to use
%%%%%%%%%%%%%%%%%%%%%%%%%%%%%%%%%%%%%%%%%%%%%%%
\usepackage{appendix}

% graphic, float package
\usepackage{graphicx}		% for setting images
\usepackage{float}			% for float objects
% \usepackage{subfloat}		% subfloat.sty 미설치 → 비활성화 (미사용)
\usepackage{subfigure}		% for adding several figures in a figure environment
\usepackage{lscape}			% for landscape type images or tables

\usepackage{enumitem}
\usepackage{xspace}			% for \model, \eg 등 매크로 trailing space 처리
\usepackage[table]{xcolor}	% tikz보다 먼저 [table] 옵션으로 로드 (옵션 충돌 방지)
\usepackage{tikz}			% 모델 구조 그림
\usetikzlibrary{arrows.meta, positioning, fit, backgrounds}

% for compact section title spacing
% \usepackage[compact]{titlesec}


% mathmetical presentation
\usepackage{gensymb}
\usepackage{amsmath}
\usepackage{amssymb}
\usepackage{amsthm}
\usepackage{exscale}
\usepackage{textcomp}		% extra symbols


% for circled number
\newcommand{\cl}[1]{\textcircled{\scriptsize #1}}


% package for using algorithmic presentation (algorithmic.sty 미설치 → 비활성화)
% \usepackage{algorithmic}
% \usepackage{algorithm}
% customize algorithmic environment
% \renewcommand{\algorithmicrequire}{\makebox[40px]{\hfill\textbf{Input :}}}
% \renewcommand{\algorithmicensure}{\makebox[40px]{\hfill\textbf{Output :}}}

% array and table presentation
\usepackage{array}
\usepackage{tabulary}
\usepackage{multirow}
% \usepackage[table]{xcolor}  % 위에서 tikz 앞에 이미 로드함 (중복 제거)
\usepackage{ctable}
\usepackage{booktabs}		% for typesetting tables at the level of publication
\usepackage{refcount}		% \getrefnumber: 표/그림 번호 발음 기반 자동 조사용